%%ha_lyman_td
Dans un premier temps, l'algorithme trie les valeurs issues du
capteur \`a point de ros\'ee (TD) pour garder
uniquement les points o\`u le capteur n'est pas en phase de recherche.
La s\'election se fait suivant
deux crit\`eres~:\\
\indent - valeur de la tension de l'\'el\'ement Peltier proche de
la valeur \`a l'\'equilibre (compris entre $V_{pelI}$ et $V_{pelS}$)\\
\indent - $Td \leq Ts$ (suppression des d\'epassements)\\
\\
La serie d'humidit\'e absolue de r\'ef\'erence est calcul\'ee \`a partir
de la s\'erie des valeurs correctes du TD.\\
A chaque pas de temps,
une r\'egression lin\'eaire est r\'ealis\'ee entre, les
N points d'humidit\'e absolue de r\'ef\'erence et, les N points
correspondants
du Lyman-$\alpha$, encadrant la valeur \`a calibrer.
La validation du r\'esultat de cette r\'egression est soumise \`a
deux crit\`eres~:\\
\indent - \'ecart-type de l'humidit\'e absolue sup\'erieur \`a 0.5 $g/m^3$
(dynamique du signal suffisante)\\
\indent - coefficient de corr\'elation sup\'erieur \`a 0.8 (qualit\'e de
la r\'egression)\\

